\subsection{06 августа. Перевал Кашкатор Северный}
\textit{Метеоусловия: утром ясно, холодно, после обеда переменная облачность, небольшой дождь, ветренно}

\begin{figure}[h!]
	\centering
	\includegraphics[angle=0, width=0.7\linewidth]{../pics/maps/06_map.jpg}
	\label{fig:06_map}
	\end{figure}

Подъём в 4:00 утра. Погода ясная, видны звезды и Венера. Холодно, палатки и трава покрыты инеем. Завтрак в 5:00, в 6:15 вышли в направлении перевала Кашкатор Северный.

\begin{figure}[h!]
	\centering
	\includegraphics[angle=0, width=0.7\linewidth]{../pics/days/06-08/IMG_8633.jpg}
	\label{fig:06_1}
	\caption{Вид на д. р. Кичи-Кызыл-Суу и место ночевки}
\end{figure}

Подъем по травянистому склону, тропа хорошо видна. Вскоре вышло солнце, стало веселее, встречали много пасущихся коров. Во время движения обсуждали местоположение группы (Остапива А.), активно пытались разглядеть ребят на склоне.

\begin{figure}[h!]
	\begin{minipage}[h!]{0.49\linewidth}
		\centering{\includegraphics[width=0.98\linewidth]{../pics/days/06-08/IMG_8640.jpg}}
	\end{minipage}
	\hfill
	\begin{minipage}[h!]{0.49\linewidth}
		\centering{\includegraphics[width=0.98\linewidth]{../pics/days/06-08/IMG_8681.jpg}}
	\end{minipage}
	\caption{Траверс на подъем травянистого склона, вид на д. р. Кичи-Кызыл-Суу}
	\label{fig:06_2}
\end{figure}

В 8:35 дошли до запасного места ночевки, отмеченного на карте как «Две палатки». Травянистый склон перешел в средне-осыпной, начали подъем под перевал Кашкатор Северный.

\begin{figure}[h!]
	\begin{minipage}[h!]{0.49\linewidth}
		\centering{\includegraphics[width=0.98\linewidth]{../pics/days/06-08/IMG_8739.jpg}}
	\end{minipage}
	\hfill
	\begin{minipage}[h!]{0.49\linewidth}
		\centering{\includegraphics[width=0.98\linewidth]{../pics/days/06-08/IMG_8755.jpg}}
	\end{minipage}
	\caption{Вид на пер. Кашкатор Северный, вид на перевал Каратакия с подхода к Кашкатору Северному}
	\label{fig:06_3}
\end{figure}



По рации связались с группой Остапива А., обменялись новостями и дальнейшими планами. Обошли по правому краю кулуара озера, начали подъем на перевальный взлет, склон мелко-осыпной.

\begin{figure}[h!]
	\centering
	\includegraphics[angle=0, width=0.7\linewidth]{../pics/days/06-08/IMG_8763.jpg}
	\label{fig:06_3-}
	\caption{«Красивые» озера на подходе к пер. Кашкатор Северный, траверс склона по мелкой сыпухе}
\end{figure}



На подъёме попадались кости животных: челюсти, позвонки. На перевальном взлёте по центру кулуара был заметен крупный скелет, вероятно, коровы. Подъем на перевал осуществлялся по правому краю кулуара, так как слева с ледника периодически летели камни.

\begin{figure}[h!]
	\begin{minipage}[h!]{0.49\linewidth}
		\centering{\includegraphics[width=0.98\linewidth]{../pics/days/06-08/IMG_8775.jpg}}
	\end{minipage}
	\hfill
	\begin{minipage}[h!]{0.49\linewidth}
		\centering{\includegraphics[width=0.98\linewidth]{../pics/days/06-08/IMG_8777.jpg}}
	\end{minipage}
	\caption{Руковод добыла позвонки, вид на пер. Каратакия и Киргизский хребет}
	\label{fig:06_4}
\end{figure}


В 11:55 достигли перевала Кашкатор Северный. Оставили записку, съели перевальную шоколадку, отдохнули.

\begin{figure}[h!]
	\begin{minipage}[h!]{0.49\linewidth}
		\centering{\includegraphics[width=0.98\linewidth]{../pics/days/06-08/IMG_8826.jpg}}
	\end{minipage}
	\hfill
	\begin{minipage}[h!]{0.49\linewidth}
		\centering{\includegraphics[width=0.98\linewidth]{../pics/days/06-08/IMG_8833.jpg}}
	\end{minipage}
	\caption{Группа на перевале Кашкатор Северный}
	\label{fig:06_5}
\end{figure}

В 12:20 начали спуск с перевала по руслу высохшей реки. В начале приятный мелкоосыпной спуск, уже к 12:55 вышли на травянистый склон в д. р. Кашка-Тёр.

\begin{figure}[h!]
	\begin{minipage}[h!]{0.49\linewidth}
		\centering{\includegraphics[width=0.98\linewidth]{../pics/days/06-08/IMG_8848.jpg}}
	\end{minipage}
	\hfill
	\begin{minipage}[h!]{0.49\linewidth}
		\centering{\includegraphics[width=0.98\linewidth]{../pics/days/06-08/IMG_8860.jpg}}
	\end{minipage}
	\caption{Спуск по мелкой осыпи, вид на д.р. Кашка-Тёр}
	\label{fig:06_6}
\end{figure}


Стало переменно-облачно, периодически капал дождик, усилился ветер, обед пропустили. Шли по травянистому склону с небольшим уклоном, встречали табуны лошадей.

\begin{figure}[h!]
	\begin{minipage}[h!]{0.49\linewidth}
		\centering{\includegraphics[width=0.98\linewidth]{../pics/days/06-08/IMG_8905.jpg}}
	\end{minipage}
	\hfill
	\begin{minipage}[h!]{0.49\linewidth}
		\centering{\includegraphics[width=0.98\linewidth]{../pics/days/06-08/IMG_8916.jpg}}
	\end{minipage}
	\caption{Спуск вдоль реки Кашка-Тёр по тропам для скота, вид на д р Джуукучак, где будет ночевка}
	\label{fig:06_7}
\end{figure}


В 15:15 дошли до крутого склона в д. р. Джуукучак – планируемое место ночевки. Спускались траверсом по хорошим тропам для скота.

\begin{figure}[h!]
	\begin{minipage}[h!]{0.49\linewidth}
		\centering{\includegraphics[width=0.98\linewidth]{../pics/days/06-08/IMG_8947.jpg}}
	\end{minipage}
	\hfill
	\begin{minipage}[h!]{0.49\linewidth}
		\centering{\includegraphics[width=0.98\linewidth]{../pics/days/06-08/IMG_8954.jpg}}
	\end{minipage}
	\caption{Спуск в д.р. Джуукучак}
	\label{fig:06_8}
\end{figure}


В 15:45 вышли к реке, ищем брод, шли вдоль реки по течению.

\begin{figure}[h!]
	\centering
	\includegraphics[angle=0, width=0.7\linewidth]{../pics/days/06-08/IMG_8957.jpg}
	\label{fig:06_9}
	\caption{Вид на д. р. Джуукучак}
\end{figure}


Наконец нашли неглубокий участок реки, к 16:30 перебрались на другой берег, встали легерем. Поужинали, легли спать рано.

\begin{figure}[h!]
	\begin{minipage}[h!]{0.49\linewidth}
		\centering{\includegraphics[width=0.98\linewidth]{../pics/days/06-08/IMG_8981.jpg}}
	\end{minipage}
	\hfill
	\begin{minipage}[h!]{0.49\linewidth}
		\centering{\includegraphics[width=0.98\linewidth]{../pics/days/06-08/IMG_9033.jpg}}
	\end{minipage}
	\caption{Брод реки Джуукучак, глубина меньше метра, лагерь в д.р.}
	\label{fig:06_10}
\end{figure}


\begin{figure}[h!]
	\centering
	\includegraphics[angle=0, width=0.7\linewidth]{../pics/days/06-08/IMG_8990.jpg}
	\label{fig:06_11}
	\caption{А коровкам брод нипочем!}
\end{figure}











\begin{tabular}{|c|c|}
	\hline
	ЧХВ: & 2:32 \\
	\hline
	ГХВ: & 3:30 \\
	\hline
\end{tabular}

\clearpage
